% Producer of the pairwise receipt (G3) of subsec:trace-rounded-comparison.
% Issue #372 (parent #56, C15). Correspondence and evidence:
% blueprint/units/YM-000026/handoff/results/C15/author/hairer-claude-372-20261114/.
\subsection{The pairwise receipt from the level budget}
\label{subsec:trace-rounded-receipt}

We now derive the receipt (G3) of
Subsection~\ref{subsec:trace-rounded-comparison} from two properties of
the funded chart witnesses of
Section~\ref{sec:ex-lattice-firesky-gallery}: a scale allowance, stated in
(A) below, and a root-test bound at each level whose size is an explicit
negative power of the clock monitor, stated in (B) below. The width clause
uses only (A) and the growth of the monitor. The two root-test bounds
\eqref{eq:trace-rounded-receipt-one}--\eqref{eq:trace-rounded-receipt-two}
use only (B) and the growth of the monitor. The selection of the chart
levels then supplies (G3) at the cost $\varepsilon_i=(i+1)^{-1}$. At the
end of the subsection we show how the replica-root estimate of
Section~\ref{sec:rank-root-normalization} yields (B), and attach its inputs to the selected funded family.

\paragraph{Charts at a fixed level.}
Throughout, $\mu_k\ge1$ is the clock monitor of (G1), with
$\mu_k\to\infty$; it is the diverging monitor $\mathcal M_k$ of
Section~\ref{sec:rank-root-normalization}. Fix a grade cutoff $d$ and
let
\[
 M=M_d=d+2
\]
be the marker allowance of the funded chart family at grade $d$: the
maximal number of markers per cell in the replica estimate of
Section~\ref{sec:rank-root-normalization}. (It is not the arity cap $l$
of Section~\ref{sec:ex-lattice-firesky-gallery}.) Consider a chart of the
family at level $k$ with base arc $b>0$. Put
\[
 m=\mu_k^{\,b},\qquad r=a_km,
\]
so that $r=g_b(k)$ in the notation of \eqref{eq:ex-lattice-firesky-gap}.
Let $a_{\ell}$ be the spacing of the middle cell of the chart. Let $w$ be
the root depth of the chart: the depth of the centered tiling on which
the root tests of its comparison receipt are taken, and at which its
chamber condition is asserted. A configuration $x\in(\mathbb R^4)^2$ is
\emph{fit for the chart} if the grid configuration $a_k\lceil mx\rceil$
is coordinatewise separated and each of its coordinates has absolute
value at most $a_\ell/8$. For the chart selected at index $i$, with
$(k,b)=(k_i,b_i)$, this is exactly fitness at level $i$ as defined in
Subsection~\ref{subsec:trace-rounded-comparison}. As there, for a fit
$x$ and $e,e'\in E_{d'}$ with $d'\le d$, we write $f_0^{N,p}$,
$f_1^{N,p}$ for the finite-volume insertions of the chart labels of $e$
and $e'$ at the two sites of $a_k\lceil mx\rceil$. We write
$j_0,j_1,\kappa$ for the reference values at the rounded configuration
$\rho_m(x)$, and $\Phi_s$ ($s\subseteq\{0,1\}$) for the parts of the
corrected polynomial defined there. Thus $\Phi_{\{c\}}=f_c^{N,p}-j_c$ and
$\Phi_{\{0,1\}}=(f_0^{N,p}-j_0)(f_1^{N,p}-j_1)-\kappa$.

We use the following two properties of the funded chart witnesses. In
both, the level threshold is fixed before the chart is chosen.
\begin{enumerate}
\item[(A)] \emph{Allowance.} There is $K_A=K_A(d)$ such that every chart
 of the family at a level $k\ge K_A$ with base arc $b$ satisfies
 \begin{equation}\label{eq:trace-receipt-allowance}
  m\,\mu_k^{2}\le\frac{a_\ell}{a_k},
  \qquad\text{equivalently}\qquad r\,\mu_k^2\le a_\ell .
 \end{equation}
\item[(B)] \emph{Level budget.} There is $K_B=K_B(d)$ with the following
 property. Let $k\ge K_B$ and consider any chart of the family at level
 $k$, any grade $d'\le d$, any basis vectors $e,e'$ of $E_{d'}$, any
 $N>k$, any period $p$, and any $x$ fit for the chart with
 $\rho_m(x)\in\mathcal W$. Then, for every nonempty
 $s\subseteq\{0,1\}$, the root test of the centered tiling at the root
 depth $w$ of the chart, for the Wilson law at $(N,p)$, satisfies
 \begin{equation}\label{eq:trace-receipt-budget}
  \|\Phi_s\|_{\mathrm{root}}\le\mu_k^{-1/(16M)} .
 \end{equation}
\end{enumerate}
The thresholds $K_A$ and $K_B$ may depend on the grade cutoff, the chosen
bases and the window $\mathcal W$. They do not depend on the chart, on
$(N,p)$, or on $x$. Since there are finitely many grades $d'\le d$ and
finitely many pairs of basis vectors, a threshold for each fixed pair
yields one for all of them by taking the maximum. This chart-uniform form
is used in the proof of Proposition~\ref{prop:trace-rounded-receipt},
where the threshold must be known before the producer is called. A
weaker form, with a threshold allowed to depend on the level curve
returned by the producer, would suffice for the first assertion of that
proposition, but not for the argument given for its second assertion.

\begin{lemma}[The width clause]\label{lem:trace-receipt-width}
Let $R>0$. Consider a chart at a level $k\ge K_A$ with base arc $b$, and
suppose $\mu_k^2\ge8R$. Let $t\in(\mathbb Z^4)^2$ be such that
$y=m^{-1}t$ is coordinatewise separated and $|y|\le R$. Then $y$ is fit
for the chart. In particular, if $\mu_k^2\ge800$, the width clause of
(G3) holds for that chart.
\end{lemma}
\begin{proof}
Since $my=t$ has integer coordinates, $\lceil my\rceil=t$, and the grid
configuration of $y$ is
\[
 z=a_k\lceil my\rceil=a_kt=a_kmy=ry .
\]
Multiplication by $r>0$ preserves the strict order of real numbers, so
$z$ is coordinatewise separated because $y$ is. Each coordinate
satisfies $|z_{c,\nu}|=r|y_{c,\nu}|\le r|y|\le rR$, because a coordinate
of a vector in $\mathbb R^{8}$ is bounded by its Euclidean norm. By
\eqref{eq:trace-receipt-allowance}, which holds because $k\ge K_A$, and
by $\mu_k^2\ge8R$,
\[
 rR\le\frac{a_\ell}{\mu_k^2}\,R\le\frac{a_\ell}{8R}\,R=\frac{a_\ell}{8}.
\]
Hence $y$ is fit. The width clause of (G3) is the case $R=100$, and
there $8R=800$.
\end{proof}

\begin{lemma}[Cost of the root tests]\label{lem:trace-receipt-cost}
Assume (B), and let $0<\varepsilon\le1$. If $k\ge K_B$ and
$\mu_k\ge\varepsilon^{-16M}$, then for every chart of the family at level
$k$ the bounds
\eqref{eq:trace-rounded-receipt-one}--\eqref{eq:trace-rounded-receipt-two},
with the root test taken at the root depth $w$ of the chart, hold with
$\varepsilon$ in place of $\varepsilon_i$. This holds for every
grade $d'\le d$, all basis vectors $e,e'$ of $E_{d'}$, all $N>k$, all
$p$, and every $x$ fit for the chart with $\rho_m(x)\in\mathcal W$.
\end{lemma}
\begin{proof}
The function $t\mapsto t^{-1/(16M)}$ is decreasing on $(0,\infty)$, so
$\mu_k\ge\varepsilon^{-16M}$ gives
$\mu_k^{-1/(16M)}\le(\varepsilon^{-16M})^{-1/(16M)}=\varepsilon$. By
(B), $\|\Phi_s\|_{\mathrm{root}}\le\varepsilon$ for each nonempty
$s\subseteq\{0,1\}$. By the identities recalled above, the left sides
of \eqref{eq:trace-rounded-receipt-one} and
\eqref{eq:trace-rounded-receipt-two} are exactly
$\|\Phi_{\{c\}}\|_{\mathrm{root}}$ and
$\|\Phi_{\{0,1\}}\|_{\mathrm{root}}$.
\end{proof}

Since $\mu_k\to\infty$, the conditions
\[
 k\ge K_A(d),\qquad k\ge K_B(d),\qquad \mu_k^2\ge800,\qquad
 \mu_k\ge\varepsilon^{-16M_d}
\]
hold for all sufficiently large $k$. We write $K(d,\varepsilon)$ for the
least $K$ such that they hold for every $k\ge K$. (The monitor is not
assumed monotone, so this is the right form of the threshold.)

\begin{proposition}[Producer of the pairwise receipt]
\label{prop:trace-rounded-receipt}
Assume (A) and (B) at every grade cutoff. Identify the root depth $w_i$
of (G4) with the root depth $w$ of the chart selected at index $i$; that
is, (G4) is required at that depth. Suppose that for every index $i$ the
gallery selection of Section~\ref{sec:ex-lattice-firesky-gallery}, with
grade cutoff $d=i$ and cost index $l=i$, takes its chart at a level
\begin{equation}\label{eq:trace-receipt-level}
 k_i\ge K\bigl(i,(i+1)^{-1}\bigr).
\end{equation}
Then the selected data satisfy (G3) at the cost
$\varepsilon_i=(i+1)^{-1}$, including its width clause. Moreover
the selection can always be made so that \eqref{eq:trace-receipt-level}
holds for every $i$, without changing the conclusions of
Lemma~\ref{lem:ex-lattice-firesky-selection}.
\end{proposition}
\begin{proof}
Fix $i$, and let $M=M_i=i+2$. The selected chart has level
$k_i\ge K(i,(i+1)^{-1})$ and base arc $b_i$, so
$m=\mu_{k_i}^{b_i}=m_i$ and the fitness of the previous paragraphs is
fitness at level $i$. By the identification in the statement, the root
test of (G3) at index $i$ is taken at the depth $w_i=w$ of the selected
chart. Since $k_i\ge K_A(i)$ and $\mu_{k_i}^2\ge800$,
Lemma~\ref{lem:trace-receipt-width} with $R=100$ gives the width clause:
every $y=m_i^{-1}t$ that is coordinatewise separated with $|y|\le100$ is
fit at level $i$. Since $k_i\ge K_B(i)$ and
$\mu_{k_i}\ge(i+1)^{16M}=\varepsilon_i^{-16M}$,
Lemma~\ref{lem:trace-receipt-cost} with $\varepsilon=\varepsilon_i$ gives
\eqref{eq:trace-rounded-receipt-one}--\eqref{eq:trace-rounded-receipt-two}
at depth $w_i$. These hold for every grade $d'\le i$, all basis vectors
of $E_{d'}$, all $N>k_i$, all $p$, and every $x$ fit at level $i$ with
$\rho_{m_i}(x)\in\mathcal W$. As $i$ was arbitrary, this is (G3).

For the second assertion we use only the output of the one-step producer
as stated in Section~\ref{sec:ex-lattice-firesky-gallery}. The chart it
returns at level $k$ and base arc $b>5$ has exact radius equal to the
clock gap $g_b(k)=a_k\mu_k^b$ of \eqref{eq:ex-lattice-firesky-gap}, and the
producer accepts every threshold $\theta>0$ for that radius. Fix $i$ and
put $K=K(i,(i+1)^{-1})$; by (A) and (B) as stated, $K$ is determined
before the producer is called. Since $\mu_k\ge1$ and $b>0$, the radius
satisfies $g_b(k)\ge a_k$. Given the threshold $\theta>0$ required by the
selection at step $i$, call the producer instead with
\[
 \theta'=\min\Bigl(\theta,\ \min_{j<K}a_j\Bigr)>0
\]
(with $\theta'=\theta$ if $K=0$); this minimum is positive because
there are finitely many $j<K$ and each $a_j>0$. If the returned level $k$
satisfied $k<K$, the returned radius would satisfy
$g_b(k)\ge a_k\ge\theta'$, contradicting $g_b(k)<\theta'$. Hence
$k\ge K$, which is \eqref{eq:trace-receipt-level}. Since
$\theta'\le\theta$, the returned output also meets the threshold $\theta$
required in the proof of Lemma~\ref{lem:ex-lattice-firesky-selection},
and every other property of the output is unchanged. Thus that lemma
applies verbatim with the reduced thresholds, and its conclusions hold
for the resulting selection.
\end{proof}

\paragraph{Supply of the level budget.}
The bound \eqref{eq:trace-receipt-budget} is the replica-root estimate of
Section~\ref{sec:rank-root-normalization} at marker allowance $M$. Write
$q_0$ for the number of replica cells, which is also the site count of
the selected output torus. With
\[
 I_k=\frac{\log\mu_k}{16M},\qquad a=e^{2I_k}=\mu_k^{1/(8M)}\ge1 ,
\]
a coefficient whose pushed marked density has the local estimate with
coefficient price $a^{q_0|I|}$ and per-site expense $X\le1+M$ satisfies
\eqref{eq:rank-root}, whose right side is
$\exp\{P_R(X)-4D\log g/q_0\}/a^{|I|}$. The harmonic term
$-4D\log g/q_0$ is nonnegative and must be paid for. That section shows
\[
 P_R(X)-4D\log g/q_0\le T_M
\]
under the following conditions, all at the output level $m$ of the clock:
the coupling bound $\beta\le4\widehat\beta_m$; the harmonic condition
$8D\widehat\beta_m/q_0\le1$, which the clock supplies eventually; the
same clip and row bound $A$ as the clock stage, with
$\theta_R\le\theta_{\rm clock}$, $\Lambda\le\Lambda_{\rm clock}$ and
$\delta\le1$; and the late-scale purchases
$\widehat\beta_m\theta_{\rm clock}^3\le1$,
$\widehat\beta_m\theta_{\rm clock}^4\le1$ and
$\Lambda_{\rm clock}\theta_{\rm clock}^2\le1$. If moreover the stage ticket
satisfies $T_M\le I_k$, then the right side of \eqref{eq:rank-root} is
at most $e^{I_k}/a^{|I|}$. For a nonempty marker subset $|I|\ge1$ and
$a\ge1$, so this is at most $e^{I_k}/a=e^{-I_k}=\mu_k^{-1/(16M)}$. The
marker subsets $I$ are the nonempty $s\subseteq\{0,1\}$; since
$M\ge2$, they are admissible. The coefficients are the one-point and
corrected two-point parts $\Phi_s$ themselves, as
Section~\ref{sec:rank-root-normalization} requires.

The coefficient price $a^{q_0|I|}$ with $X\le1+M$ on the same pushed
density is Proposition~\ref{prop:allarity-output-price}, with $u=M$. It
is obtained from the scalar envelope \eqref{eq:allarity-density-envelope}
at the marker weight \eqref{eq:allarity-marker-weight}. That weight is
used from the launch at the middle level down to the output torus. The
leading interval above the middle level is run at a smaller weight. The
proposition's inputs, each on the common history of the chart witness,
are listed in the status paragraph at the end of
Section~\ref{sec:allarity-marked-propagation}.

\begin{proposition}[The receipt on the selected chart family]
\label{prop:trace-selected-rounded-receipt}
Use the common history and the funded panels of
Proposition~\ref{prop:trace-joint-chart-purchase}, with the ticket requests
\eqref{eq:rank-root-ticket-requests} included before the history is
chosen. For each fixed grade cutoff $d$, restrict to its three panels
and their offsets $t\in[0,1]$. Then (A) and (B) hold, with thresholds
uniform over those charts, all $N>k$, all periods, and all the fit
configurations in the stated window. The root test in (B) is at the
same depth $w$ as the chamber statement
Lemma~\ref{lem:ex-lattice-firesky-chamber}.
Consequently the selected gallery can satisfy (G3) at every index,
without changing its background history or extraction.
\end{proposition}
\begin{proof}
Fix $d$ and put $M=d+2$. Choose the constants of the leading
recurrence before its coupling threshold: after fixing its multiplier
and small-weight exponent $U$, take
$q\ge\max\{6,\kappa+4,\kappa+Un+2\}$, where $\kappa$ is the
upper step-count exponent and $n=2$ below. These are the initial,
spare-weight and launch margins of
Proposition~\ref{prop:allarity-launch}. The arcs of these panels have the fixed upper
bound $\bar b=\max_{0\le h\le2}(b_{d,h}+1)$. All applications below use
this bound; no uniformity over unbounded arcs is asserted. Put
$s=\widehat\beta_k$ and $H=\mu_k$. In the notation of
Section~\ref{sec:allarity-marked-propagation}, take $n=2$ marks per
replica cell and $u=M$. Thus $n\le u$, and its large marker weight is
$\omega=H^{1/(8M)}=a$. Its output level is the depth $w$ defined in
Subsection~\ref{subsec:ex-lattice-firesky-witness}, not the middle
level $\ell$. The output torus has exactly $q_0$ sites, one per cell
of this centered tiling.

\emph{The marked density and the five propagation inputs.}
We check the inputs listed at the end of
Section~\ref{sec:allarity-marked-propagation} on this particular
history. Use the corrected labels solving
\eqref{eq:ex-lattice-firesky-solves}, their actual insertions, and their
signed replica copies throughout.
\begin{enumerate}
\item The original mixed-grade marked expansion and its census payment
 in Section~\ref{sec:ex-lattice-firesky-gallery} give the terminal
 multi-marker representation of these insertions, including the
 ordinary and forced pieces. Its payment is
 \eqref{eq:trace-original-marked-payment}; the endpoint change of
 weights is \eqref{eq:trace-endpoint-marked-budget}. Both apply to
 the same supported terms, and retain the scalar and unmarked density.
 For the initial regular remainder, use the endpoint error estimate
 there with its parameter $g=q$ and its reserve order at least $q+4$:
 $H^{-D-8}A_0+d_s\le H^{-q}$. The forced payment requires a separate argument. Perform the
 endpoint regrouping before summing the marked factors, in the finite
 squarefree marker algebra with its weighted $\ell^1$ norm. This norm
 is submultiplicative, so the connected-set majorant and the gas
 bijection in Proposition~\ref{prop:restart-endpoint} apply unchanged.
 Every nonzero forced component contains a selected flag factor $X_z$:
 an old forced inventory contains an old flag, hence a new flag, and
 a mask-error factor vanishes unless its full support contains a new
 flag. Write $u_k=a_{\rm coast}\mathrm{hot}(s)^2$. The endpoint proof
 gives $|X_z|\le e^{-2u_k}$; reserve $e^{-u_k}$ in each component.
 The remaining site-factor price is precisely the $b_s$ of the
 gallery endpoint estimate. The old-support, marked-root and
 inherited-factor majorants are its $d_s,t_s,m_s$, evaluated in the
 same weighted norm. They do not have to be divided by $G_k$: their
 component already has the reserved flag factor. With the endpoint's
 two additional error terms their combined residual majorant is
 \[
  r_s=2H^{-3}+t_s+m_s+b_s\le H^{-1},\qquad
  Z(B)r_s\le1,\qquad e^{r_s}-1\le1.
 \]
 The connected sum is therefore at most
 $e^{-u_k}(e^{r_s}-1)\le e^{-u_k}=G_k$.
 This is a bound for the entire marked forced gas, not only its
 zero-marker part. At zero markers the original endpoint identity
 and scalar normalization are unchanged. The small weight
 $\omega_-\le1$ can only reduce the already paid marker norms.
 The terminal substitution of
 Proposition~\ref{prop:rank-selected-tested-readout} replaces only the
 near singleton by its masked raw polynomial and its two paid errors;
 the wide, forced and higher-mode terms are retained. It supplies the
 common polynomial and low-grade reading, rather than independently
 chosen maps. Its stamp and the calibrated polynomial price of
 Section~\ref{sec:allarity-coordinate-calibration} give the terminal
 source purchase and the principal power bound used by
 Lemma~\ref{lem:allarity-principal-power}.

 At every leading rung the actual compiler application is the one in
 the ``Room and the actual restarted steps'' argument of
 Section~\ref{sec:ex-lattice-firesky-gallery}. It keeps the path,
 clip, tube, forest, tag and positive loading. In particular
 Proposition~\ref{prop:restart-actual-closure} supplies the incoming
 neutral law and the genuine remainder bound; the joint-warning
 transfer is Lemma~\ref{lem:actual-neutral-warning}. These are the
 hypotheses of Section~\ref{sec:actual-local-compiler-handoff},
 not assumptions about an arbitrary coarse representation.

 For the local discrepancy use the guarded-law alternative
 Corollary~\ref{cor:allarity-guarded-purchase}, not exact transport
 by a normalized cutoff integral. The room and adjacent-profile
 comparison of Section~\ref{sec:ex-lattice-firesky-gallery} verify
 the hypotheses of Lemma~\ref{lem:allarity-pair-reader}. In particular
 the seed-free store, pose and enlarged guard rooms are those of
 Section~\ref{sec:actual-reader-dawn-selection}, transported with
 the actual paths, forest and tag. That lemma supplies one joint
 Gaussian reader array, its mean forecast and short covariance,
 with coordinate and covariance errors at most $H^{-f}$ and guard
 loss at most $H^{-2f}$, uniformly on the healthy coarse fields.

 Here is the component identification needed for this application.
 In the actual mixed compiler an entirely leading component has no
 body letter, and compatibility excludes background packets. Its
 integral is a product of the guarded principal values on the common
 Haar--Gaussian sample space. On a positive guard the field is
 healthy on the entire prescribed mask support, so that value is
 $p_I(\xi_{t,I})$. Whole-mode selection sums these products over
 partitions, with the same nonnegative forest/tag weights for every
 product, including the empty product, whose integral is one.
 These are exactly $M_S$ of Lemma~\ref{lem:allarity-guarded-gaussian}.
 Gauge invariance removes the patch gauge from each physical value.
 Finite Wick contraction with the same forecast and short covariance
 gives its $M^*_S$, the prescribed next polynomial moments.
 This identifies only the principal component, not the full
 interacting block density as Gaussian. The common loading is the
 one already fixed by the compiler; it is not replaced by a separate
 interacting partition function.

 The calibrated coefficient and covariance bounds give $A,B\le H^c$
 for a fixed $c$ on this finite menu. If $r$ is the degree of the
 polynomial in the guarded lemma, its connected error is
 $y\le C H^{cr-f}$. Fix the required remainder exponent $q$ first
 and choose $f>q+cr$. Then $2^nR_{e_*}H^qy\le1$ eventually.
 This is the local discrepancy purchase at the reserved support
 exponent. Realignment to the prescribed next polynomial propagates
 this error in the regular and forced remainders; no second history
 of normalized cutoff operators is iterated.
 Reader separation is Lemma~\ref{lem:allarity-halo-separation}.
\item The same reader comparison and the exact low-grade reading just
 identified give the moment error in
 Section~\ref{sec:allarity-cumulant-comparison}. The actual per-step
 compiler and terminal stamp give the healthy local error in
 Section~\ref{sec:marked-healthy-cumulant-launch}. The polynomial
 error estimates in Section~\ref{sec:ex-lattice-firesky-gallery}
 allow their fixed losses to be paid before the accuracy exponent is
 chosen, so both errors are at most $H^{-3/2}$ at the middle level.
 These are precisely the two hypotheses of
 Lemma~\ref{lem:allarity-corrected-pieces}; they concern the same
 principal polynomial and reference cumulant as in item 1.
\item Choose the support reserve before making the local compiler
 purchases: fix $0<\sigma\le1$, set $e_*=e+1+\sigma$ and $E=e+2$,
 and take $\alpha\ge8T_E$. Apply
 Proposition~\ref{prop:allarity-compiler-reserve} to the actual steps
 of item 1, with local support exponent $E+8$. The finite-stage
 block and entry purchases of
 Proposition~\ref{prop:trace-joint-chart-purchase} include these
 constants for each of its adjacent stage pairs. Thus the upper
 calculation and the terminal purchase hold at $e_*$ on the same
 letters. This verifies (E); Lemma~\ref{lem:allarity-spare-run} supplies the
 enlarged-exponent run, and
 Lemma~\ref{lem:allarity-remainder-restored} restores the marker
 weight. Coarse footprints contain the projected fine supports by
 the actual mixed compiler; Lemma~\ref{lem:allarity-reader-containment}
 gives the required containment. Halo separation supplies (S), so
 Proposition~\ref{prop:allarity-launch} applies.
\item The same reserved compiler supplies the lower calculation at
 exponent $e$ on $[w,\ell)$. Its remainder recurrence is
 Lemma~\ref{lem:allarity-lower-propagation}; its step count and
 forced inputs are supplied by Lemmas~\ref{lem:allarity-lower-count-free}
 and~\ref{lem:allarity-lower-forced}. The constant coupling, stage
 changes and normalizing scalar are those of
 Proposition~\ref{prop:restart-actual-closure}.
\item Corollary~\ref{cor:rank-trace-all-levels} supplies the lawful
 representations on that same course, at every level of $[w,k]$.
 At $w$ its clip and row bound are exactly those of $\pi_{J(w)}$,
 and it satisfies the coupling, threshold, loop and additive bounds
 required by Proposition~\ref{prop:rank-root-output-ticket}.
\end{enumerate}
These applications are uniform over the stated window: it is bounded
and its two points have distance at least $10^{-2}$. The finite-reader
and census estimates use these bounds, not a positive lower bound for
each coordinate difference. The latter differences need only be nonzero
for the chamber and ordering conditions. There are finitely many basis
pairs and grades through $d$. All the corresponding constants and all
three adjacent-stage cases are included in the finite-stage maxima
before the common history is chosen. The joint purchase then gives a
level tail independent of the chart offset, $N,p$ and $x$.

The propagated density is the push of
$\exp_*(-D_R)\prod_c(1+f_c^{N,p}z_c)$ times the actual background density,
with $D_R$ the reference cumulant polynomial. The field-independent
factor $\exp_*(-D_R)$ commutes with every push. The complete marked
push identities used in item 1 preserve every mode, not just the
singletons, and the unmarked coefficient preserves the partition
function. By Lemma~\ref{lem:allarity-invoice-parts} the three
nonempty coefficients are exactly the $\Phi_s$ of this subsection.
Proposition~\ref{prop:allarity-output-price} therefore gives, on this
lawful representation at $w$, the coefficient price $a^{q_0|s|}$ and
expense $X\le1+n\le1+M$.

\emph{Clock, ticket and uniform root depth.}
Lemma~\ref{lem:ex-lattice-firesky-clock-fields} supplies (A), the
plateau and the two-stage bracket uniformly for $b\le\bar b$.
The bracket gives $J(w)\ge J(k)-1$. In particular, for any fixed
level $m_0$, once $J(k)-1>J(m_0)$ every such chart has $w>m_0$.
This proves the uniform version of $w\to\infty$ needed for the
late-scale and harmonic purchases; one does not choose a different
tail for each chart curve. The ticket proof in
Proposition~\ref{prop:rank-root-output-ticket} is uniform in $w\le k$:
$J(w)\le J(k+1)$ lets the same countable payment buy
$T_M(\pi_{J(w)})\le I_k$. Combining it with the lawful representation
just obtained gives
\[
 \|\Phi_s\|_{\rm root}
 \le \frac{\exp\{P_{P_w}(X)-4D\log g/q_0\}}{a^{|s|}}
 \le \frac{e^{I_k}}{e^{2I_k|s|}}
 \le e^{-I_k}=\mu_k^{-1/(16M)}
 \qquad(s\ne\varnothing).
\]
This proves (B), including its finite-volume uniformity. Finally the
chamber lemma uses exactly the centered tiling at $w$. Its mesh
threshold is allowed in (G4); it does not weaken the root estimate,
which holds for every $N>k$. Proposition~\ref{prop:trace-rounded-receipt}
now supplies (G3), by decreasing the requested radius thresholds before
each chart selection. Only the selected charts change, not the history
or its extraction.
\end{proof}

\paragraph{Status.}
The conditional producer Proposition~\ref{prop:trace-rounded-receipt}
continues to apply to any family satisfying (A) and (B).
Proposition~\ref{prop:trace-selected-rounded-receipt} verifies those
properties on the selected funded family and identifies its root depth.
Thus (G3) is supplied for that family's rounded comparison. This
application retains the canonical selected-history and representation
constructions; it does not assert the receipt for an arbitrary chart
with merely the same numerical scales.
